% !TEX program = pdflatex
\documentclass[preprint,12pt]{elsarticle}

\usepackage{amssymb}
\usepackage{amsmath}
\usepackage{booktabs}
\usepackage{multirow}
\usepackage{mathptmx}           % Times New Roman font
\usepackage{float}              % [H] placement for tables/figures at end
\usepackage[hidelinks]{hyperref}% remove coloured hyperlinks
\usepackage{xcolor}
\usepackage{setspace}
\usepackage[explicit]{titlesec}
\usepackage[top=2.54cm, bottom=2.54cm, left=2.54cm, right=2.54cm]{geometry}
\usepackage{textcomp}

%% Remove section numbering
\setcounter{secnumdepth}{0}

%% Section: bold, centred (template style)
\titleformat{\section}[block]{\normalfont\large\bfseries\centering}{}{0em}{#1}
\titlespacing*{\section}{0pt}{18pt}{6pt}

%% Subsection: italic, left-aligned (template style)
\titleformat{\subsection}[block]{\normalfont\normalsize\itshape}{}{0em}{#1}
\titlespacing*{\subsection}{0pt}{12pt}{4pt}

%% Subsubsection: italic, left-aligned
\titleformat{\subsubsection}[block]{\normalfont\normalsize\itshape}{}{0em}{#1}
\titlespacing*{\subsubsection}{0pt}{8pt}{2pt}

%% Double spacing throughout (manuscript submission style)
\doublespacing

%% Patch pprintMaketitle (preprint mode): remove hrules and author block
%% (authors and abstract are typeset manually after \end{frontmatter})
\makeatletter
\long\def\pprintMaketitle{\clearpage
  \iflongmktitle\if@twocolumn\let\columnwidth=\textwidth\fi\fi
  \resetTitleCounters
  \def\baselinestretch{1}%
  \printFirstPageNotes
  \begin{\elsarticletitlealign}%
    \thispagestyle{pprintTitle}%
    \def\baselinestretch{1}%
    \Large\@title\par\vskip18pt
  \end{\elsarticletitlealign}%
  \gdef\thefootnote{\arabic{footnote}}%
}
\makeatother

%% Point to the figure directory
\graphicspath{{/home/josh/NTU/114-2/MLHMA/final_proj/}}

\journal{Final Project of Machine Learning in Human Motion Analysis}

\begin{document}

\begin{frontmatter}

\title{Is Ground Reaction Force Alone Sufficient for Clinical Balance Assessment?
       A Subject-Independent Benchmark of Deep Learning Architectures for
       COM-COP Relative Motion Prediction}

\end{frontmatter}

%%-------------------------------------------------------------
%% Cover page: authors, affiliations, word count, date, email
%%-------------------------------------------------------------
\begin{center}
b10404022 \quad Hung-Lin Chang$^{1}$ \\[6pt]
b12508010 \quad Yu-I Lin$^{2}$ \\[6pt]
b12508016 \quad Yu-Chieh Ho$^{2}$ \\[6pt]
b12508026 \quad Wei-Hsuan Tai$^{2}$
\end{center}

\vspace{1em}
\begin{center}
{\small\itshape
$^{1}$Department of Clinical Laboratory Sciences and Medical Biotechnology, Taipei, Taiwan \\[4pt]
$^{2}$Department of Biomedical Engineering, Taipei, Taiwan}
\end{center}

\vspace{3em}
\begin{center}
Word Count: $\sim$6{,}000
\end{center}

\vspace{2em}
\begin{center}
June 21, 2026
\end{center}

\vspace{2em}
\begin{center}
E-mail: \{b10404022, b12508010, b12508016, b12508026\}@ntu.edu.tw
\end{center}

%%-------------------------------------------------------------
%% Abstract page (manual, so \clearpage works)
%%-------------------------------------------------------------
\clearpage
\begin{center}
  {\normalfont\large\bfseries Abstract}
\end{center}

\noindent\textit{Background:}
Accurate assessment of postural balance requires measuring the displacement of the
centre of mass (COM) relative to the centre of pressure (COP), yet the optical
motion capture systems conventionally used for this purpose are expensive,
laboratory-bound, and operationally demanding.
Ground reaction force (GRF) signals from a single force plate represent a widely
available alternative, but whether they are sufficient for motion-capture-free
balance assessment under subject-independent conditions has not been systematically
established.

\bigskip
\noindent\textit{Research Question:}
This study addresses four specific questions: (1)~Can GRF signals alone predict
COM-COP relative displacement within a clinically useful accuracy range under strict
subject-independent evaluation?
(2)~Which deep learning architectural family performs best, and are performance
differences statistically significant?
(3)~Can real-time causal inference match the offline upper bound without meaningful
accuracy loss?
(4)~How many additional subjects are required to reach a clinical accuracy milestone?

\bigskip
\noindent\textit{Methods:}
Eight established deep learning architectures and three real-time causal adaptation
strategies were benchmarked on 106 walking trials from 12 healthy subjects under
strict leave-one-subject-out cross-validation (12 folds), ensuring no subject
information leaks into model selection or evaluation.

\bigskip
\noindent\textit{Results:}
The best offline model (BiLSTM with attention) achieved 8.68~mm mediolateral RMSE
and Pearson $r = 0.990$ under subject-independent evaluation.
Five of six pairwise comparisons among the top four architectures were not
statistically significant (Wilcoxon $p = 0.20$ to $0.97$), indicating the task's
information ceiling at current data scale.
CausalCNN was not significantly different from BiLSTM\_Attn in the
anterior-posterior direction ($p = 0.677$) at zero latency with only 75K
parameters, though it remains significantly behind the AP-best offline model
CNN1D ($p = 0.005$).
Power-law extrapolation indicates approximately 61 subjects may be needed for
BiLSTM\_Attn to reach 8~mm mediolateral RMSE.

\bigskip
\noindent\textit{Significance:}
GRF signals alone are sufficient for motion-capture-free dynamic balance assessment.
Architecture choice is a secondary concern at current data scale; the primary
bottleneck is dataset size.
These findings provide a concrete roadmap and recruitment target for clinical
translation.

\bigskip
\noindent\textit{Keywords:} ground reaction force; centre of mass; centre of
pressure; deep learning; gait analysis; balance assessment; real-time prediction

%%-------------------------------------------------------------
%% Highlights page
%%-------------------------------------------------------------
\clearpage
\begin{center}
  {\normalfont\large\bfseries Highlights}
\end{center}
\renewcommand{\labelitemi}{\textbullet}
\begin{itemize}
  \item GRF signals alone achieve sub-9~mm mediolateral RMSE for COM-COP prediction under strict subject-independent evaluation
  \item Five of six pairwise comparisons among the top four architectures are not statistically significant ($p = 0.20$ to $0.97$), indicating the information ceiling at current data scale
  \item CausalCNN is not significantly different from BiLSTM\_Attn (the ML-best offline model) in the anterior-posterior direction ($p = 0.677$) at zero latency with only 75K parameters, though it remains significantly behind the AP-best offline model CNN1D ($p = 0.005$)
  \item Power-law extrapolation of summary means provides an indicative recruitment target of approximately 61 subjects for clinical-grade accuracy
\end{itemize}
\clearpage

%%=============================================================
\section{Introduction}
\label{sec:intro}
%%=============================================================

\subsection{Background and Motivation}

In the fields of human biomechanics and neural control, accurate assessment of
postural balance and locomotor stability is fundamental for fall prevention in
older adults, diagnosis of neurological disorders, and optimization of
rehabilitation training.
However, obtaining precise dynamic balance measurements in clinical practice
remains technically challenging.
Currently, the gold standard for measuring the trajectory of the body's center of
mass (COM) is the Optical Motion Capture System.
Although this technique provides high measurement accuracy, its clinical
application is limited by several factors, including high equipment costs, strict
laboratory space requirements, and the time-consuming procedures involved in
placing reflective markers and performing subsequent data processing.
Furthermore, optical motion capture systems generally confine subjects to
controlled laboratory environments, making it difficult to capture natural
movements during daily activities or in home settings.
These limitations hinder the widespread implementation of balance assessment
technologies in clinical and real-world applications~\cite{s23104933}.

To overcome the limitations of conventional measurement systems, understanding
the dynamic relationship between the center of mass (COM) and the center of
pressure (COP) has become a major focus in biomechanical studies of postural
control.
During quiet standing, the human body is commonly modeled as an inverted pendulum.
Within this framework, the COM represents the center of the body's mass
distribution, and its position and velocity can be regarded as important state
variables describing postural stability.
In contrast, the COP is defined as the effective point of application of the
ground reaction force (GRF) on the support surface and can be considered an
external manifestation of the neural control strategy resulting from neuromuscular
regulation of ankle joint moments.
Under conditions of quiet standing or small-amplitude body sway, the COM
trajectory is typically smooth and dominated by low-frequency components because
of the body's inertia, whereas the COP exhibits larger amplitudes and
higher-frequency fluctuations, reflecting the nervous system's rapid adjustments
to postural deviations~\cite{s23104933,labrozzi2024center}.

The primary objective of this study is to investigate the dynamic interaction
between COM and COP, whose physical significance lies in quantifying the
restorative torque required to maintain postural stability.
Under the small-angle approximation of the inverted pendulum model, the restorative
torque can be expressed as $\tau = mg(\mathrm{COM} - \mathrm{COP})$, where the
horizontal distance between COM and COP serves as the moment arm through which
the ground reaction force generates a restoring effect.
When the COM deviates from its equilibrium position, the central nervous system
modulates neuromuscular activity to adjust the COP location, thereby producing
restorative torque that regulates the acceleration and direction of COM motion and
helps maintain dynamic stability~\cite{labrozzi2024center}.

Therefore, the relative motion between COM and COP reflects the dynamic process
by which the neural control system maintains balance through adjustments in plantar
pressure distribution.
By integrating GRF measurements with inverted pendulum dynamics, it may be
possible to accurately estimate COM-COP relative motion using only a limited set
of anthropometric parameters, thereby reducing reliance on large-scale optical
motion capture systems.
Compared with approaches based on multi-segment body models and anthropometric
regression tables, physics-based estimation methods may also mitigate errors
arising from inter-individual differences in mass distribution.
Consequently, such methods have the potential to provide a more cost-effective and
real-time solution for dynamic balance assessment in clinical
settings~\cite{s23104933,labrozzi2024center}.

\subsection{Research Gap}

Despite the clinical appeal of GRF-based balance assessment, prior work has not
addressed this question with the rigour required for clinical translation.
Existing studies either rely on within-subject evaluation (which conflates
memorisation of individual gait patterns with generalisation), compare a single
architecture without systematic ablation, or do not isolate the performance gap
between offline and real-time deployment.
This study addresses all three gaps in a unified experimental framework.

\subsection{Research Questions}

This study addresses the question through a systematic benchmark with four
specific research questions:

\begin{description}
  \item[RQ1 (Feasibility)] Can triaxial GRF signals alone predict COM-COP
        relative displacement within a clinically useful accuracy range under
        strict subject-independent evaluation?
  \item[RQ2 (Architecture)] Which deep learning architectural family performs
        best, and is the performance difference between architectures statistically
        significant or masked by inter-subject variability?
  \item[RQ3 (Real-Time Deployment)] Can non-causal offline models be adapted
        to zero-latency causal inference without significant accuracy loss, and
        what drives performance in the causal setting?
  \item[RQ4 (Data Requirements)] How many additional subjects would be needed
        to reach a clinically meaningful accuracy milestone?
\end{description}

\subsection{Summary of Findings}

Our analysis of 12 subjects (106 trials) under leave-one-subject-out (LOSO)
cross-validation yields the following answers:

\begin{description}
  \item[RQ1] \textbf{Yes.} The best offline model achieves 8.68~mm ML RMSE and
        Pearson $r = 0.990$, confirming that GRF alone is sufficient for
        motion-capture-free balance assessment.
  \item[RQ2] Five of six pairwise comparisons among the top four architectures
        (BiLSTM\_Attn, CNN1D, BiLSTM, TCN\_LSTM) are \textbf{not statistically
        significant} ($p = 0.20$ to $0.97$); the one significant pair (CNN1D vs
        TCN\_LSTM, $p = 0.034$) reflects CNN1D's dilated-receptive-field advantage.
        Architecture choice is a secondary concern until more data become available.
  \item[RQ3] CausalCNN achieves AP performance \textbf{not significantly different}
        from BiLSTM\_Attn ($p = 0.677$) at zero latency with only 75K parameters,
        though it remains significantly worse than CNN1D (the AP-best offline model,
        $p = 0.005$). Architecture design, not offline knowledge transfer, is the
        decisive factor.
  \item[RQ4] Power-law extrapolation of summary means indicates ${\sim}61$ subjects
        may be needed for BiLSTM\_Attn to reach 8~mm ML RMSE; this serves as an
        indicative recruitment target subject to the uncertainty of extrapolation
        well beyond the observed range.
\end{description}

%%=============================================================
\section{Related Work}
\label{sec:related}
%%=============================================================

\subsection{Physics-Based COM Estimation}

The classical approach to recovering COM kinematics from a force plate is to
double-integrate the GRF vector using Newton's second law
($\mathbf{F} = m\ddot{\mathbf{r}}_\mathrm{COM}$).
This approach is straightforward in principle but requires force moments
($M_x$, $M_y$) to compute the COP ($\mathrm{COP}_x = -M_y / F_z$); without
them, the COP cannot be separated from the COM trajectory.
Standard force plates do report moments, but insole-based systems typically do not,
making physics-based COP estimation infeasible in portable settings.
Furthermore, double integration accumulates drift: small constant errors in
measured forces grow quadratically over time, making the method unreliable even
over a single gait cycle~\cite{labrozzi2024center}.
The baseline comparison below confirms errors exceeding 60~mm when moments are unavailable.

\subsection{Data-Driven Balance and Gait Analysis}

Machine learning has been applied to balance and gait analysis across a range of
sensor modalities~\cite{s23104933,labrozzi2024center}.
Recurrent architectures (LSTM, GRU) are the dominant choice for sequential
biomechanical signals, with demonstrated utility in COM trajectory
estimation~\cite{labrozzi2024center}, gait phase detection, and fall risk
scoring~\cite{s23104933}.
Convolutional approaches based on temporal convolutional networks (TCN) have shown
competitive performance on similar sequence-to-sequence biosignal
tasks~\cite{bai2018empirical}, with the advantage of parallelisable training.
Transformers have more recently been applied to gait and biosignal modelling,
but their data requirements are larger by a considerable margin~\cite{vaswani2017attention},
and their advantage over recurrent models on small clinical datasets ($n < 50$)
has not been consistently demonstrated, motivating the recurrent-first architecture
choices evaluated in this study.

\subsection{Gaps Motivating This Benchmark}

Despite this activity, three gaps remain unaddressed in the existing literature:

\begin{enumerate}
  \item \textbf{Subject-independent evaluation.}
        Most GRF-based COM/COP prediction studies report within-subject or
        random-split validation, which conflates memorisation of individual gait
        patterns with generalisation.
        Leave-one-subject-out evaluation, required for clinical deployment, has not
        been systematically applied in this task.
  \item \textbf{Architecture comparison under identical conditions.}
        No prior work compares more than two or three architectures on the same
        COM-COP prediction task with identical preprocessing, training, and
        evaluation protocols, making cross-study comparisons unreliable.
  \item \textbf{Real-time deployment analysis.}
        The gap between offline (non-causal) accuracy and zero-latency causal
        deployment has not been quantified, nor have the relative contributions of
        architectural design versus offline knowledge transfer been isolated.
\end{enumerate}

This study addresses all three gaps in a unified experimental framework.

%%=============================================================
\section{Materials and Methods}
\label{sec:methods}
%%=============================================================

\subsection{Participants and Data Acquisition}
\label{subsec:data}

Gait data were obtained from a dataset of 13 healthy adults who performed level
walking on a single force plate.
Each trial recorded triaxial GRF together with the corresponding COM and COP
trajectories for one gait cycle.
Each subject's data file also included the recording sampling frequency (200 or
240~Hz, varying by subject; see Table~\ref{tab:dataset}) and body weight, which
were used to time-normalise the gait cycle and to scale the GRF signals,
respectively.
All trials were resampled to a uniform 100-timestep representation to eliminate
the effect of sampling rate differences.

Of the 13 subjects, one was excluded because one of the recorded trials contained
NaN values; since the cause of the missing data could not be confirmed and the
remaining trials for that subject could not be independently verified, all trials
from that subject were removed.
The final dataset comprised \textbf{12 subjects}, each contributing 8 to 9 trials,
for a total of \textbf{106 trials}.

Detailed demographic information beyond body weight (i.e., age, sex, and
height) was not provided with the dataset.
Body height was instead estimated from the mean vertical COM position during gait,
based on the anthropometric relationship reported by Dempster~\cite{dempster1955},
and used only to compute a height-normalised error metric for evaluating
subject-independent prediction accuracy.

\subsection{Signal Processing and Preprocessing}
\label{subsec:preproc}

Raw \texttt{.mat} files contained the following fields (Table~\ref{tab:fields}).
The following preprocessing steps were applied:
\begin{enumerate}
  \item \textbf{GRF normalisation}: $\mathbf{g} = \mathbf{GRF} / (m \cdot g_0)$,
        where $m$ is body mass (kg) and $g_0 = 9.81$~m\,s$^{-2}$, removing
        inter-subject force scaling.
  \item \textbf{COM-COP displacement}: $\boldsymbol{\delta} =
        (\mathbf{COM}_{xy} - \mathbf{COP}) / 1000$, converting mm to metres.
  \item \textbf{Temporal normalisation}: each trial was resampled from its native
        length (159 to 263 frames) to exactly 100 timesteps via linear interpolation,
        aligning all trials to 0 to 100\% of the gait cycle regardless of sampling
        rate.
  \item Any trial containing NaN values caused the entire subject to be excluded
        at load time.
\end{enumerate}

Table~\ref{tab:dataset} summarises the final dataset statistics after exclusion.

\subsection{Model Architectures}
\label{subsec:archs}

Eight architectures were evaluated with shared base hyperparameters
(\texttt{hidden\_size}~=~64, \texttt{num\_layers}~=~2, \texttt{dropout}~=~0.2).
All models map a normalised GRF sequence of shape $(T=100, 3)$ to a COM-COP
displacement sequence of shape $(T=100, 2)$.
Table~\ref{tab:archs} provides an overview.

\noindent\textbf{ANN} applies a two-hidden-layer MLP independently at each
timestep, serving as a lower-bound baseline with no temporal context.

\noindent\textbf{RNN and LSTM} are unidirectional recurrent networks.
LSTM's input, forget, and output gates enable long-range dependency modelling that
vanilla RNN lacks.

\noindent\textbf{BiLSTM} runs forward and backward passes simultaneously; each
timestep can attend to the full gait cycle, rendering the model non-causal.

\noindent\textbf{BiLSTM\_Attn} adds a multi-head self-attention layer~\cite{vaswani2017attention}
(4 heads, embed dim 128) with residual connection and LayerNorm on top of BiLSTM,
allowing the model to learn which timesteps are most informative:
\begin{equation}
  \mathbf{h} = \mathrm{LayerNorm}\!\left(\mathbf{h}_\mathrm{BiLSTM}
              + \mathrm{MHA}(\mathbf{h}_\mathrm{BiLSTM})\right).
\end{equation}

\noindent\textbf{Transformer} uses Pre-LN encoder layers~\cite{vaswani2017attention}
(2 layers, 4 heads, FFN dim 256) preceded by a linear input projection and
sinusoidal positional encoding.

\noindent\textbf{CNN1D} is a six-block dilated temporal convolutional network
(TCN)~\cite{bai2018empirical} with dilation rates $\{1, 2, 4, 8, 16, 32\}$ and
Pre-LN residual connections, yielding a total receptive field of 253 timesteps
(2.5$\times$ the gait cycle length).

\noindent\textbf{TCN\_LSTM} stacks four TCN blocks (local feature extraction)
followed by a BiLSTM (global temporal integration), combining complementary
inductive biases in a single architecture.

\subsection{Real-Time Adaptation Strategies}
\label{subsec:realtime}

The four best offline models are all non-causal, requiring the complete gait cycle
before any prediction.
Three strategies adapt this knowledge to zero-latency causal deployment:

\begin{enumerate}
  \item \textbf{CausalCNN}: CNN1D with \emph{symmetric} padding replaced by
        \emph{left-only} (causal) padding; all other weights are identical to the
        offline CNN1D.
        Each timestep sees only past and current GRF, with a receptive field of
        253 steps (75K params).
  \item \textbf{CausalAttn}: Unidirectional LSTM backbone (hidden~=~128) with a
        causally-masked self-attention layer.
        The hidden dimension is doubled relative to each direction of the teacher's
        BiLSTM (64 per direction) so that the output width stays 128 and the
        attention/norm/fc weights can be transferred directly from the offline
        BiLSTM\_Attn teacher (201K params); the LSTM weights are randomly initialised
        because the shapes differ.
        The wider single-direction LSTM actually has \emph{more} parameters than the
        two narrower halves of the teacher's BiLSTM, giving CausalAttn ${\sim}267$K
        parameters in total.
  \item \textbf{KD\_LSTM (h128)}: A standard LSTM (hidden~=~128) trained with
        knowledge distillation~\cite{hinton2015distilling} using a composite loss:
        \begin{equation}
          \mathcal{L} = (1-\alpha)\,\mathrm{MSE}(\hat{y}, y)
                       + \alpha\,\mathrm{MSE}(\hat{y},\,\hat{y}_\mathrm{teacher}),
                       \quad \alpha = 0.5,
        \end{equation}
        where $\hat{y}_\mathrm{teacher}$ are soft labels from the offline
        BiLSTM\_Attn (200K params).
\end{enumerate}

\subsection{Training Protocol}
\label{subsec:training}

All models were optimised with AdamW (lr~=~0.001, weight decay~=~$10^{-4}$) and
mean-squared-error (MSE) loss.
Gradient norms were clipped to 1.0.
Learning rate followed a cosine annealing schedule from 0.001 to $10^{-5}$ over a
maximum of 200 epochs, with early stopping (patience~=~20 on internal validation
MSE).

\noindent\textbf{Normalisation.}
Both inputs and targets were independently z-scored per fold, per channel, using
statistics computed from the training set only.
Target normalisation is critical: the AP range is approximately $3\times$ the ML
range, and without it AP gradients dominate training.
BiLSTM\_Attn improved from 11.74~mm (worst) to 8.68~mm (best) in ML RMSE after
target normalisation was introduced.

\noindent\textbf{Data augmentation} (applied only to the training set, $4\times$
original volume):
\begin{itemize}
  \item \textbf{Left-to-right flip}: negate $F_x$ and $\delta_x$ simultaneously to
        simulate mirror-symmetric gait.
  \item \textbf{Gaussian noise}: $\sigma = 0.01$ in normalised space, simulating
        sensor measurement error.
  \item \textbf{Time warping}: $\pm 10\%$ stretch/compress via linear
        interpolation, increasing temporal diversity.
\end{itemize}
MixUp regularisation ($\lambda \sim \mathrm{Beta}(0.4, 0.4)$, applied with
probability 0.5 per batch) further reduced overfitting.

\subsection{Cross-Validation and Evaluation Metrics}
\label{subsec:eval}

All models were evaluated under \textbf{leave-one-subject-out cross-validation
(LOSO, 12 folds)}: each subject serves as the held-out test set exactly once,
with the remaining 11 subjects used for training (80\%) and internal validation
(20\%).
The test set was never seen during training or model selection.

Statistical significance between models was assessed using the Wilcoxon
signed-rank test across the 12 per-fold RMSE values.

The following metrics were computed on each test fold:
\begin{itemize}
  \item \textbf{RMSE X / Y (mm)}: root mean squared error in the ML and AP
        directions.
  \item \textbf{Pearson $r$ X / Y}: linear correlation between predicted and
        ground-truth trajectories.
  \item \textbf{Angular error ($^\circ$)}: mean absolute angular deviation of the
        predicted 2-D displacement vector from the ground-truth vector.
  \item \textbf{Height-normalised RMSE (\% height)}: RMSE / estimated body
        height $\times 100$, a dimensionless cross-subject metric.
\end{itemize}

%%=============================================================
\section{Results}
\label{sec:results}
%%=============================================================

\subsection{Non-DL Baselines}
\label{subsec:baselines}

To contextualise the need for deep learning, three non-DL methods were evaluated
under the same LOSO protocol as a separate preliminary analysis
(Table~\ref{tab:baselines}).

Physics-based double-integration fails (X~=~62.9~mm) because the input contains
only GRF magnitudes, not moments ($M_x$, $M_y$).
COP recovery requires force moments (i.e., $\mathrm{COP}_x = -M_y / F_z$);
without them, any model must implicitly learn the COP trajectory from GRF temporal
patterns, a task that is beyond the scope of physics-only methods.
Pointwise linear regression captures a strong linear component (17.5~mm), but the
remaining gap to the best DL models ($\sim$8.7~mm) represents the non-linear,
temporally global gait structure that only deep sequence models can learn.

\subsection{Offline Deep Learning Models}
\label{subsec:offline}

Table~\ref{tab:offline} summarises LOSO performance across all eight offline
architectures, ordered by ML RMSE.

The top four models (BiLSTM\_Attn, CNN1D, BiLSTM, TCN\_LSTM, ordered by ML RMSE)
show five of six pairwise ML comparisons that are not statistically significant
(Wilcoxon $p = 0.20$ to $0.97$), indicating that the information ceiling of the
current dataset has been approached.
The one significant pair is CNN1D versus TCN\_LSTM ($p = 0.034$), reflecting
CNN1D's dilated-receptive-field advantage in the ML direction.
CNN1D also leads convincingly in the AP direction (15.21~mm versus 17.15~mm for
BiLSTM\_Attn; $p = 0.002$), demonstrating that its large receptive field captures
the propulsion-braking dynamics that unfold over longer temporal windows.
Performance across all models and subjects is further illustrated in
Figures~\ref{fig:offline} and~\ref{fig:heatmap}.

\subsection{Real-Time Adaptation}
\label{subsec:rt_results}

Table~\ref{tab:realtime} compares the three causal adaptation strategies against
the LSTM baseline (best causal offline model) and the offline upper bound in ML
(BiLSTM\_Attn, 8.68~mm).
Note that for the AP direction CNN1D (15.21~mm) outperforms BiLSTM\_Attn
(17.15~mm), making CNN1D the stricter AP reference.

\noindent\textbf{CausalCNN achieves the best AP performance among causal models}
(16.81~mm), which is not significantly different from BiLSTM\_Attn's AP
($p = 0.677$) but remains significantly worse than CNN1D's AP ($p = 0.005$),
reflecting the irreducible information gap between causal and offline inference.
CausalAttn leads in the ML direction at 9.55~mm, a margin of only 0.87~mm below
the offline best.
A combined comparison of all offline and real-time models is shown in
Figure~\ref{fig:allmodels}.

\subsection{Data Scaling Analysis}
\label{subsec:scaling}

To quantify how additional data would benefit each architecture, we subsampled the
training pool to $k \in \{1, 2, 3, 5, 7, 9, 11\}$ subjects (3 repeats each) and
fitted power-law learning curves $\mathrm{RMSE} = a \cdot k^{-b} + c$ to the
per-$k$ summary means (Table~\ref{tab:scaling}; Figure~\ref{fig:scaling}).
Extrapolated values beyond $k = 11$ are inherently uncertain and should be
interpreted as indicative recruitment targets rather than precise predictions.

Most models have not plateaued at $k = 11$, confirming that additional subjects
would improve performance.
CNN1D is the exception: it saturates near $k = 2$ (its dilated inductive bias
captures local patterns from few subjects but reaches a cross-subject ceiling of
${\sim}9.1$~mm early; k=7 to 11: 9.43 to 9.08~mm, floor extrapolation 9.12~mm).
BiLSTM\_Attn shows the largest remaining headroom (2.1~mm gap to its floor) and
the highest long-term scaling potential.

%%=============================================================
\section{Discussion}
\label{sec:discussion}
%%=============================================================

\subsection{Answer to RQ1: Clinical Feasibility}

The benchmark establishes that GRF alone achieves sub-9~mm ML RMSE and
$r \geq 0.988$ under strict subject-independent evaluation.
For context, the typical COM-COP displacement in the ML direction is
approximately 52~mm, measured as the standard deviation of ML displacement across
all trials in this dataset; a prediction error of 8.7~mm corresponds
to 17\% of this typical amplitude, which is within the accuracy typically accepted
for clinical gait analysis instruments~\cite{s23104933}.
This result provides a clear answer to RQ1: \textbf{GRF-only balance assessment is clinically
feasible for laboratory walking in healthy adults, and optical motion capture may be
replaceable for this specific task under these conditions.}

\subsection{Answer to RQ2: Architecture Choice}
\label{subsec:disc_offline}

Five of six pairwise comparisons among the top four models ($p = 0.20$ to $0.97$)
are not statistically significant, indicating that for most pairs the expected
accuracy difference is smaller than the trial-to-trial variability across subjects.
\textbf{Architecture choice is therefore a secondary concern at this data scale;
the primary bottleneck is dataset size.}
The one significant pair (CNN1D versus TCN\_LSTM, $p = 0.034$) reflects CNN1D's
dilated-receptive-field advantage rather than a general architectural hierarchy.

CNN1D is the strongest architecture in the AP direction: its mean AP RMSE
(15.21~mm) is 1.94~mm lower than BiLSTM\_Attn's (17.15~mm, $p = 0.002$),
and it also leads TCN\_LSTM ($p = 0.339$, not significant) and BiLSTM
($p = 0.052$, borderline).
The 253-step receptive field of CNN1D spans 2.5 gait cycles, enabling it to
capture the full propulsion-braking dynamics that recurrent models with limited
effective context cannot match.

One training design choice had larger impact than any architecture selection:
\textbf{target normalisation.}
Because the AP range is approximately $3\times$ larger than the ML range,
unnormalised training allows AP gradients to dominate, collapsing ML accuracy.
BiLSTM\_Attn went from 11.74~mm (worst among all models) to 8.68~mm (best in ML)
after per-channel z-score normalisation of targets, a result that underscores how
training protocol choices can outweigh architecture selection on small clinical
datasets.

\subsection{Answer to RQ3: Real-Time Deployment}
\label{subsec:disc_rt}

\textbf{Architecture design dominates over knowledge transfer.}
Ablation experiments isolate the contribution of offline model knowledge by
comparing paired models with and without transfer:
CausalAttn with weight transfer versus trained from scratch: $\Delta = 0.13$~mm
($p = 0.733$, not significant);
KD\_LSTM h128 versus LSTM h128 trained from scratch: $\Delta = 0.56$~mm
($p = 0.064$, not significant).
In both cases, switching from the LSTM baseline architecture to a richer causal
architecture accounts for 94\% and 76\% of the total gain, respectively.

This result can be understood from an information-theoretic perspective.
KD trains the student on teacher soft labels computed from the \emph{full} GRF
sequence, yet at test time the causal student has access only to GRF$[0{:}t]$.
There is an irreducible information gap:
\begin{equation}
  I\!\left(\mathbf{g}_{t{+}1:T};\;\boldsymbol{\delta}_t \mid \mathbf{g}_{0:t}\right) > 0,
\end{equation}
meaning the teacher's predictions depend on information the student can never
access.
The marginal 0.56~mm benefit from KD arises from label smoothing (softer targets
reduce overconfidence), not from genuine knowledge transfer of temporal context.

\textbf{KD requires sufficient student capacity.}
h64 (51K params) yielded zero improvement ($p = 0.970$) over the LSTM baseline;
h128 (200K params) yielded 8\% improvement; h256 (794K params) provided
negligible additional gain.
Capacity saturates near hidden~=~128, and the KD benefit is marginal regardless
of capacity.

\subsection{Answer to RQ4: Data Requirements}

The power-law learning curves confirm that most models have not plateaued at
11 training subjects; CNN1D is the exception, approaching saturation near $k=2$.
BiLSTM\_Attn, the architecture with the deepest extrapolated floor (7.01~mm),
is estimated to require approximately 61 subjects to reach 8~mm ML RMSE,
based on power-law extrapolation of the per-$k$ summary means.
This extrapolation is sensitive to the small number of data points and should be
treated as an indicative recruitment target rather than a precise threshold.
CNN1D saturates near $k = 2$ due to its strong local inductive bias; it is the
architecture of choice when data are scarce, but its cross-subject ceiling of
${\sim}9.1$~mm limits long-term scalability.
\textbf{Future data collection efforts should therefore prioritise BiLSTM\_Attn
as the primary model, targeting 30 subjects as an intermediate milestone
(predicted ${\sim}8.4$~mm ML RMSE).}

\subsection{Clinical Applications}
\label{subsec:applications}

\subsubsection{Rehabilitation Tracking and Fall Risk Assessment}
\label{subsubsec:rehab}

Based on the framework proposed in this study, the primary application of
predicting COM-COP relative motion from Ground Reaction Forces (GRF) is
longitudinal rehabilitation tracking for Total Knee Arthroplasty (TKA) patients.
Effective post-operative recovery strictly relies on ``restoring bilateral
symmetry'' as a core indicator~\cite{tka2025}.
By utilizing our prediction model, clinicians can efficiently evaluate patients at
fixed post-surgery intervals during gait analysis.
The trajectory mismatch in COM-COP relative motion between the healthy and
affected sides can serve as a highly sensitive, objective symmetry indicator, with
the ultimate therapeutic goal of minimizing this discrepancy.
Crucially, rather than benchmarked against a generic healthy database, which often
introduces confounding variables such as lifelong gait habits and subtle anatomical
variations, rehabilitation progress is evaluated using the patient's own
contralateral side as a personalized baseline.

The second proposed application of this predictive framework is fall risk
assessment in older adults.
Age-related degradation in dynamic balance control significantly elevates fall
risks within this population.
Although conventional clinical assessment tools such as the Berg Balance Scale
and the Timed Up and Go (TUG) test are widely utilized, they rely predominantly
on subjective scoring and time duration, failing to capture underlying, real-time
biomechanical data.
To bridge this clinical gap, the methodology developed in this study utilizes
predicted Ground Reaction Forces (GRF) to calculate advanced dynamic stability
metrics, specifically moving from the Center of Mass (COM) to the Extrapolated
Center of Mass (XCoM), defined as $\mathrm{XCoM} = \mathrm{COM} + \dot{\mathrm{COM}} / \sqrt{g/l}$,
and ultimately determining the Margin of Stability (MoS).
By calculating XCoM (which accounts for both COM position and velocity) and
quantifying its distance from the Base of Support boundary (BoS), this approach
offers an objective, high-precision, and quantitative diagnostic tool for proactive
fall risk screening in geriatric care.

\subsubsection{Real-Time Device Control and Veterinary Settings}
\label{subsubsec:device}

A third application direction is real-time device control.
Plantar-pressure smart insoles already provide a practical means of acquiring
GRF-like signals outside the laboratory~\cite{s22249725}, and the zero-latency
causal models evaluated in this study, which process only past GRF
($\mathbf{g}_{0:t}$) at each timestep, demonstrate that the prediction task is
feasible without access to future measurements, suggesting the architecture could
operate under similar real-time constraints.
This is relevant for applications such as exoskeletons and powered orthoses, where
maintaining the user's COM within a stable region is a recognised balance control
objective~\cite{10.3389/fnbot.2021.751642}, and for powered lower-limb prostheses, which
already rely on real-time GRF thresholds to switch between control
states~\cite{10342504}.

A fourth application, suggested during course discussion, extends this approach
beyond human gait to zoo and veterinary settings.
Force plates are already an established tool for quantifying limb loading and
detecting lameness in animals such as horses, dogs, and
elephants~\cite{veterinary2022}, but full marker-based motion capture is often
impractical in these settings, since markers cannot be reliably attached to large,
free-ranging, or non-cooperative animals.
A model that estimates COM-COP relative motion directly from GRF, without
requiring marker tracking, could therefore support lameness and welfare monitoring
in zoo environments where a force plate can be embedded into an existing walkway.
Applying the present approach to non-human species would require retraining on
species-specific gait data, as body proportions, limb number, and locomotor
patterns differ substantially from human bipedal walking.

These applications would require further validation before clinical or commercial
use.
The present dataset was limited to 12 healthy adults walking under controlled
laboratory conditions, and the model has not been tested on wearable hardware, on
populations with gait impairment, or on tasks beyond level walking.
Extending the training data to include pathological gait patterns and validating
the model against pressure-insole hardware are necessary next steps before this
approach can be considered for the applications discussed above.

\subsection{Limitations}
\label{subsec:limits}

\textbf{Small and homogeneous sample.}
The dataset comprises 12 healthy adults performing controlled level walking, which
is a narrow sample relative to the diversity of clinical populations.
Data scaling analysis confirms that most models continue to improve at $k = 11$
(CNN1D approaches saturation near $k = 2$), and BiLSTM\_Attn is projected to
reach ${\sim}8.4$~mm with ${\sim}30$ subjects.
Generalisability to older adults, patients with gait pathology, or individuals
with substantial weight variation remains to be established.

\textbf{Causal versus non-causal models.}
The four best offline models are non-causal, requiring the complete gait cycle
before prediction.
For post-trial clinical gait analysis (e.g., TKA rehabilitation assessment or
laboratory-based fall risk screening), this is acceptable since data are processed
after each recording.
For streaming applications (e.g., exoskeleton control or wearable monitoring),
CausalCNN and CausalAttn are the recommended deployment choices based on the
results in Table~\ref{tab:realtime}.

%%=============================================================
\section{Conclusions}
\label{sec:conclusion}
%%=============================================================

To our knowledge, we present one of the first subject-independent benchmarks
for predicting COM-COP relative motion from GRF signals alone.
Answering the four research questions in turn:
\textbf{(RQ1)} GRF alone achieves sub-9~mm ML RMSE ($r \geq 0.988$) under
strict leave-one-subject-out (LOSO) evaluation, confirming that motion-capture-free
dynamic balance assessment is clinically feasible.
\textbf{(RQ2)} Five of six pairwise comparisons among the top four offline
architectures are not statistically significant; dataset size, not architecture
design, is the current bottleneck.
\textbf{(RQ3)} CausalCNN achieves AP performance not significantly different from
BiLSTM\_Attn ($p = 0.677$) at zero latency with only 75K parameters, though it
remains significantly behind the AP-best offline model CNN1D ($p = 0.005$);
ablation experiments confirm that architectural design outweighs offline knowledge
transfer as the performance driver for causal models.
\textbf{(RQ4)} Power-law extrapolation indicates that approximately 61 subjects may
be needed for BiLSTM\_Attn to reach the 8~mm clinical milestone, with 30 subjects
as an indicative intermediate target (${\sim}8.4$~mm); these figures are subject to
the uncertainty of extrapolation well beyond the observed range.

Together, these findings define a clear research roadmap: the GRF-to-COM-COP
prediction is approaching the ceiling of what current data permit, and the
next step is expanding the dataset (e.g., to pathological gait and older adults)
rather than exploring new architectures.

%%=============================================================
\bibliographystyle{elsarticle-num}
\bibliography{refs}
%%=============================================================

%%=============================================================
%% TABLES
%%=============================================================
\clearpage
\begin{center}
{\normalfont\large\bfseries TABLES}
\end{center}

\begin{table}[H]
\centering
\caption{Raw data fields recorded per trial.}
\label{tab:fields}
\begin{tabular}{llll}
\toprule
Field & Shape & Unit & Description \\
\midrule
\texttt{GRF} & $(T, 3)$ & N   & Triaxial ground reaction force: $F_x$ (ML), $F_y$ (AP), $F_z$ (vertical) \\
\texttt{COM} & $(T, 3)$ & mm  & Centre of mass position: X (ML), Y (AP), Z (vertical) \\
\texttt{COP} & $(T, 2)$ & mm  & Centre of pressure: X (ML), Y (AP) \\
\texttt{BodyWeight} & $(1,1)$ & kg & Subject body mass \\
\texttt{frame\_rate} & $(1,1)$ & Hz & Sampling rate (200 or 240) \\
\bottomrule
\end{tabular}
\end{table}

\begin{table}[H]
\centering
\caption{Final dataset statistics after exclusion.}
\label{tab:dataset}
\begin{tabular}{ll}
\toprule
Property & Value \\
\midrule
Subjects (original / final) & 13 / 12 \\
Trials per subject & 8 to 9 \\
Total trials & 106 \\
Sampling rate & 240~Hz (subjects 1 to 2); 200~Hz (subjects 3 to 13, original IDs) \\
Excluded subject & Subject 10 (NaN values; all trials removed; original IDs retained) \\
Raw sequence length & 159 to 263 frames per trial \\
Normalised sequence length & 100 timesteps \\
Input channels & 3 ($F_x$, $F_y$, $F_z$) \\
Output channels & 2 (ML and AP COM-COP displacement, m) \\
\bottomrule
\end{tabular}
\end{table}

\begin{table}[H]
\centering
\caption{Overview of the eight offline deep learning architectures.}
\label{tab:archs}
\begin{tabular}{llrl}
\toprule
Model & Temporal modelling & Params & Causal \\
\midrule
ANN          & Pointwise MLP (no context)              &   4,546 & Yes \\
RNN          & Vanilla recurrent (no gates)            &  12,866 & Yes \\
LSTM         & Gated recurrent (3 gates)               &  51,074 & Yes \\
BiLSTM       & Bidirectional LSTM                      & 134,914 & No  \\
BiLSTM\_Attn & BiLSTM + multi-head self-attention      & 201,218 & No  \\
Transformer  & Full-sequence self-attention (Pre-LN)   & 100,482 & No  \\
CNN1D        & 6-block dilated TCN (RF = 253 steps)    &  75,394 & No  \\
TCN\_LSTM    & 4-block TCN $\rightarrow$ BiLSTM        &  75,522 & No  \\
\bottomrule
\end{tabular}
\end{table}

\begin{table}[H]
\centering
\caption{Comparison of non-DL baselines and DL models (mean $\pm$ std, 12 folds).}
\label{tab:baselines}
\begin{tabular}{lcc}
\toprule
Method & RMSE X (mm) & RMSE Y (mm) \\
\midrule
Physics: double-integrate GRF   & $62.9 \pm 12.0$ & $271.5 \pm 15.4$ \\
Pointwise linear regression     & $17.5 \pm 2.9$  & $40.5 \pm 5.0$   \\
Linear regression + $\pm 2$-step window & $13.7 \pm 2.7$ & $29.7 \pm 4.3$ \\
ANN (DL, no temporal context)   & $13.0 \pm 2.5$  & $32.7 \pm 4.4$   \\
LSTM (DL, causal)               & $10.8 \pm 2.2$  & $18.3 \pm 5.0$   \\
\textbf{BiLSTM\_Attn (DL, full context)} & $\mathbf{8.7 \pm 2.2}$ & $17.2 \pm 4.7$ \\
\textbf{CNN1D (DL, full context)}        & $8.7 \pm 2.7$ & $\mathbf{15.2 \pm 4.2}$ \\
\bottomrule
\end{tabular}
\end{table}

\begin{table}[H]
\centering
\caption{LOSO performance of offline models (mean $\pm$ std, 12 folds).
         Bold indicates the best value per column.}
\label{tab:offline}
\begin{tabular}{lcccccc}
\toprule
Model & RMSE X (mm) & RMSE Y (mm) & $r_X$ & $r_Y$ & Angle ($^\circ$) & norm X (\%) \\
\midrule
\textbf{BiLSTM\_Attn}
      & $\mathbf{8.68 \pm 2.16}$ & $17.15 \pm 4.74$
      & $\mathbf{0.990}$ & $0.976$ & $9.63$ & $\mathbf{0.520}$ \\
CNN1D & $8.73 \pm 2.69$ & $\mathbf{15.21 \pm 4.20}$
      & $0.989$ & $\mathbf{0.981}$ & $\mathbf{8.53}$ & $0.525$ \\
BiLSTM & $9.22 \pm 2.37$ & $16.34 \pm 4.76$
      & $0.988$ & $0.978$ & $8.90$ & $0.553$ \\
TCN\_LSTM
      & $9.30 \pm 2.47$ & $15.87 \pm 4.26$
      & $0.988$ & $0.979$ & $9.19$ & $0.559$ \\
Transformer & $9.77 \pm 2.17$ & $18.51 \pm 4.54$
      & $0.985$ & $0.973$ & $10.30$ & $0.586$ \\
LSTM  & $10.81 \pm 2.21$ & $18.34 \pm 5.02$
      & $0.982$ & $0.971$ & $10.18$ & $0.647$ \\
RNN   & $11.19 \pm 2.34$ & $18.81 \pm 4.43$
      & $0.982$ & $0.970$ & $10.22$ & $0.670$ \\
ANN   & $13.01 \pm 2.47$ & $32.70 \pm 4.41$
      & $0.975$ & $0.907$ & $18.44$ & $0.774$ \\
\bottomrule
\end{tabular}
\end{table}

\begin{table}[H]
\centering
\caption{Real-time adaptation results (mean $\pm$ std, 12 folds).
         $p$-values from Wilcoxon signed-rank test vs LSTM baseline and
         vs offline upper bound (BiLSTM\_Attn) in the ML direction (X).
         $\ast$~$p<0.05$, $\ast\ast$~$p<0.01$, $\ast\ast\ast$~$p<0.001$.}
\label{tab:realtime}
\begin{tabular}{lccccc}
\toprule
Method & RMSE X (mm) & RMSE Y (mm) & $r_X$ & vs LSTM & vs Offline \\
\midrule
BiLSTM\_Attn (offline UB) & $8.68 \pm 2.16$ & $17.15 \pm 4.74$ & $0.990$ & \textemdash & \textemdash \\
\midrule
\textbf{CausalAttn} & $\mathbf{9.55 \pm 2.44}$ & $18.46 \pm 4.87$ & $\mathbf{0.986}$
                    & $p<0.001^{***}$ & $p=0.021^{*}$ \\
\textbf{CausalCNN}  & $9.57 \pm 2.31$ & $\mathbf{16.81 \pm 4.27}$ & $0.987$
                    & $p<0.001^{***}$ & $p=0.007^{**}$ \\
KD\_LSTM h128       & $9.79 \pm 2.09$ & $17.56 \pm 4.37$ & $0.986$
                    & $p=0.001^{**}$  & $p=0.002^{**}$ \\
LSTM (causal baseline) & $10.81 \pm 2.21$ & $18.34 \pm 5.02$ & $0.982$ & \textemdash & \textemdash \\
\bottomrule
\end{tabular}
\end{table}

\begin{table}[H]
\centering
\caption{ML RMSE learning curve (mean $\pm$ std over 12 folds $\times$ 3 repeats)
         and power-law extrapolation.
         Representative values shown; the full sweep used $k \in \{1,2,3,5,7,9,11\}$.}
\label{tab:scaling}
\begin{tabular}{lcccc}
\toprule
$k$ (training subjects) & LSTM & BiLSTM\_Attn & CNN1D & \\
\midrule
1  & $17.43 \pm 3.34$ & $13.66 \pm 3.98$ & $12.75 \pm 3.66$ & \\
3  & $14.44 \pm 2.47$ & $11.18 \pm 2.62$ & $10.18 \pm 2.66$ & \\
7  & $13.06 \pm 2.44$ & $9.78  \pm 2.30$ & $9.43  \pm 2.75$ & \\
\textbf{11 (current)} & $\mathbf{11.94 \pm 2.34}$ & $\mathbf{9.11 \pm 2.13}$
                      & $\mathbf{9.08 \pm 2.69}$ & \\
\midrule
Extrapolated floor    & 8.42~mm & \textbf{7.01~mm} & 9.12~mm & \\
Subjects to reach 8.0~mm & unreachable & ${\sim}61$ & unreachable & \\
\bottomrule
\end{tabular}
\end{table}

%%=============================================================
%% FIGURE CAPTIONS
%%=============================================================
\clearpage
\begin{center}
{\normalfont\large\bfseries FIGURE CAPTIONS}
\end{center}

\noindent Figure~1: LOSO performance of the eight offline models across all
evaluation metrics.  Error bars indicate $\pm 1$ standard deviation over 12 folds.

\bigskip
\noindent Figure~2: Per-subject RMSE heatmap for all eight offline models.
Each cell corresponds to one LOSO test fold.

\bigskip
\noindent Figure~3: Combined comparison of all offline and real-time models.
The dashed line marks the offline upper bound (BiLSTM\_Attn).
CausalCNN's AP bar is not significantly different from the upper bound ($p = 0.677$).

\bigskip
\noindent Figure~4: Learning curves for LSTM, BiLSTM\_Attn, and CNN1D as a
function of the number of training subjects.  Shaded areas indicate $\pm 1$
standard deviation; dashed lines show power-law extrapolations.

%%=============================================================
%% FIGURES
%%=============================================================
\clearpage
\begin{figure}[H]
\centering
\includegraphics[width=\linewidth]{compare_results.png}
\caption{LOSO performance of the eight offline models across all evaluation
         metrics. Error bars indicate $\pm 1$ standard deviation over 12 folds.}
\label{fig:offline}
\end{figure}

\clearpage
\begin{figure}[H]
\centering
\includegraphics[width=\linewidth]{compare_subject_heatmap.png}
\caption{Per-subject RMSE heatmap for all eight offline models. Each cell
         corresponds to one LOSO test fold.}
\label{fig:heatmap}
\end{figure}

\clearpage
\begin{figure}[H]
\centering
\includegraphics[width=\linewidth]{compare_all_methods.png}
\caption{Combined comparison of all offline and real-time models. The dashed line
         marks the offline upper bound (BiLSTM\_Attn). CausalCNN's AP bar is
         not significantly different from the offline upper bound (BiLSTM\_Attn, $p = 0.677$).}
\label{fig:allmodels}
\end{figure}

\clearpage
\begin{figure}[H]
\centering
\includegraphics[width=\linewidth]{data_scaling.png}
\caption{Learning curves for LSTM, BiLSTM\_Attn, and CNN1D as a function of
         the number of training subjects. Shaded areas indicate $\pm 1$ standard
         deviation; dashed lines show power-law extrapolations.}
\label{fig:scaling}
\end{figure}

\end{document}
